\structure{ИССЛЕДОВАТЕЛЬСКАЯ~ЧАСТЬ}

В исследовательской части рассматриваются теоретические основы разрабатываемой
системы распределённых вычислений и методики верификации соответствия TLA+
спецификации и исходного кода. Сначала описывается модель разрабатываемой
системы распределённых вычислений, включающая узлы, каналы связи, допустимые
отказы и условия поддержания согласованного состояния. Затем проводится обзор и
сравнительный анализ существующих алгоритмов консенсуса, обеспечивающих
согласованность данных и состояний между узлами распределённой системы.

Далее рассматриваются существующие подходы к верификации соответствия
формальной TLA+ спецификации и программной реализации. После этого
формулируется математическая постановка задачи верификации соответствия в
терминах системы переходов, трассы исполнения программы и критерия
согласованности поведения реализации со спецификацией. Таким образом,
исследовательская часть задаёт теоретическую основу как для проектирования
системы распределённых вычислений, так и для последующей разработки средств
верификации её соответствия формальной модели.
